\begin{annexitem}{annex:p12.har.poly-word-order}{thm:p12-har-poly-word-order}{The finite support, exact word programs and every order case}
\begin{proof}
W, R and P have only the displayed nonrecursive fields or direct
recursive fields. The positive source branch generator consequently
gives exactly the three recursor signatures in the body. For example
\(\mathrm{Rec}_W\) z f(i::w) contracts to f i w(\(\mathrm{Rec}_W\) z f w); no recursive
value is supplied for the index i. \(\mathrm{Rec}_R\) applies its one branch to
c,w; \(\mathrm{Rec}_P\) z f(r\#p) applies f to r,p and its one recursive value.
Prop induction has these same original arguments with the respective
P(w) or P(p) hypothesis instead of a data value. This constructs
the finite support without importing map/append/sorting operations.
\(I_C\) has seed v:C and constant step returning u:C, so its asserted
code-zero/code-one computations and typing are actual nat-recursion
instances, including C=W,P,T. Empty and field projections are those
recursors with their recursive argument ignored.

For complete word-code terms let
\(z_e:=\mathrm{Rec}_W\mathsf{one}(\lambda j:\mathsf{nat}.
\lambda v:W.\lambda r:\mathsf{nat}.O)\) and
\(z_l:=\mathrm{Rec}_W O(\lambda j:\mathsf{nat}.
\lambda v:W.\lambda r:\mathsf{nat}.\mathsf{one})\).
Both have type W-to-nat. Weq is \(\mathrm{Rec}_W\) at that target with seed \(z_e\)
and step
\[
 \lambda i:\mathsf{nat}.\lambda u:W.\lambda q:W\to\mathsf{nat}.
 \lambda v:W.\mathrm{Rec}_W O
  (\lambda j:\mathsf{nat}.\lambda v':W.\lambda z:\mathsf{nat}.
       \mathsf{andb}(\mathsf{eqb}\ i\ j)(q v'))v.
\]
Wlt has seed \(z_l\) and the identical typed step except that its last
expression is orb(ltb i j)(andb(eqb i j)(q v')). The inner targets
are nat and z is ignored. Thus the only outer recursive call is q v'
at the smaller first word; the displayed four cases of each function
follow by actual constructor computations. Word induction on the first
word with all second words quantified gives the two ranges, using
the complete boolean ranges in every cons/cons branch.

To supply the equality premise rather than a word-equality oracle,
define length by \(\mathrm{Rec}_W\) O(lambda i,w,r.S r), head by its case with
empty default O and cons result i, and tail by its case with empty
default [] and cons result w. These are typed data functions. Equality
of [] with i::w would give O=S(length w) by congruence, impossible
by the proved natural constructor separation. Equality of i::u with
j::v gives i=j and u=v by head/tail congruence. Conversely these
two equations give equality of the conses by two constructor congruences
and transitivity. Word induction and the four displayed cases now prove
weq=one iff word equality: empty/empty is reflexive, mixed shapes are
impossible, and cons/cons uses andb, eqb and the induction hypothesis
with precisely those two constructor implications. Its zero/negative
equivalence follows from the proved range and separation of O/one,
by the same two-branch Prop argument used for the earlier equality code.

For lex order the exact endpoint case propositions are false for
[]/[], true for []/cons, false for cons/[], and
\(i<j\vee(i=j\wedge u<_{\rm lex}v)\) for cons/cons.
They follow by induction on the declared relation proof: its nil,
head-less and equal-head/tail-less branches give respectively trivial,
the left injection and the right pair. A data case predicate on the
two endpoints is the motive; constructor separation handles mismatched
shapes. Conversely the four endpoint cases rebuild the corresponding
listed constructor, with FF elimination in the impossible cases.
Hence word induction, the nat ltb/eqb specifications, and orb/andb
specifications prove wlt=one iff this actual relation. Its zero/negative
equivalence again follows from range, not DN.

Irreflexivity is word induction: the empty case has no constructor;
the cons case would give either i<i or the smaller tail's reflexive
inequality. Both contradict the respective induction/ordinal result.
For transitivity induct on the first word with both other words
quantified. If it is empty, a first inequality forces the middle word
to be a cons; the second then forces the last to be a cons, and the
nil constructor proves the conclusion. If all three are conses with
heads i,j,k, their two inequalities give four cases. Strict/strict
uses nat strict transitivity i<j<k. Strict/equal and equal/strict
transport their strict head inequality along the given head equation.
Equal/equal composes the two head equations and uses the smaller-tail
induction hypothesis for the two tail inequalities. A middle or last
empty word cannot arise in this first-cons case. This proves every
transitivity branch. Asymmetry follows from transitivity/irreflexivity.
For trichotomy, induct on the first word and case the second. Empty
cases give equality or the nil inequality. With two conses use nat
head trichotomy; distinct heads give the corresponding head inequality.
Equal heads use tail trichotomy: tail equality gives cons equality,
and either tail inequality gives its equal-head constructor. All
equation transports have Prop targets. This is the full strict total
order used by insertion, independent of word sortedness.
\end{proof}
\end{annexitem}

\begin{annexitem}{annex:p12.har.poly-coefficients}{thm:p12-har-poly-coefficients}{The actual interpretation typing and coefficient induction}
\begin{proof}
Chi's seed is 0A:A and its step lambda k:nat.lambda a:A.a+1A
has type nat-to-A-to-A. V's seed is 1A and its word branch applies
the given nu:nat-to-A and multiplication to its recursive A value.
The row branch applies chi and V then multiplication, so \(\mathrm{Rec}_R\)
returns A. L's zero seed and row-value-plus-recursive-value branch
likewise have the exact \(\mathrm{Rec}_P\) target A. This supplies every declared
interpretation, with no claimed homomorphism among its typing premises.

Chi O=0A is a constructor computation. Chi one=0A+1A=1A uses the
additive identity. To prove chi(c+d)=chi c+chi d, induct on c
with d quantified. At O the two sides are chi d and 0A+chi d,
equal by identity. At S c the left, by the nat plus clause and
chi computation, is chi(c+d)+1A. Its induction rewrite is
(chi c+chi d)+1A; addition associativity and commutativity of the
last two summands give (chi c+1A)+chi d, the right side.
This is a finite three-summand rearrangement using the actual axioms.

For chi(cd)=chi c times chi d, induct on c with d quantified.
At O use chi0=0A and left zero annihilation. At S c, source times
computes cd+d. The proved chi addition equation and induction
hypothesis give (chi c times chi d)+chi d. Unit and distributivity
rewrite this as (chi c+1A)times chi d, which is the asserted
successor result. Thus the multiplicative clause is derived, not
postulated. In the Nat instance chi0=O; if chi c=c, its successor
is c+one=S c by the proved source addition clause. Nat induction
therefore gives chi n=n in primitive equality. All equation rewrites
are congruence/transitivity/Prop transport, never equality reflection.
\end{proof}
\end{annexitem}

\begin{annexitem}{annex:p12.har.poly-word-merge}{thm:p12-har-poly-word-merge}{Both nested merge recursions, sortedness and actual interpreted multiplication}
\begin{proof}
In the body term, outer seed lambda v.v has type W-to-W. At its
cons step, q:W-to-W is the smaller-first-word recursion. The inner
seed i::u:W and branch lambda j,v',z.\(I_W\)(leb i j)(i::q(j::v'))(j::z)
have the required types W and nat-to-W-to-W-to-W. Thus its inner
result is W and abstraction over v gives W-to-W. All displayed
merge cases follow from these actual nested recursors; in the false
branch z is the inner smaller-second-word result, not an unlicensed
same-size recursive call.

Element membership is generated by the head occurrence and the tail
occurrence, both in Prop. Nested word induction (first word, and
second word inside its cons step) proves the membership equivalence.
The empty cases return the other word. In the true branch a head i
is retained and the smaller-first recursion merges its tail with the
whole second word, so occurrences are exactly that head, first tail
or second word. In the false branch a head j is retained and the
smaller-second recursion merges the whole first word with the second
tail, giving exactly the same union. Both directions introduce or
eliminate the corresponding head/tail constructor; no set/list library
is used and duplicates are not removed.

From \(\mathrm{Sorted}_W\)(i::u), every occurrence in u is at least i. Prove this
by word induction on u: no occurrence at []; at j::v the adjacent
bound i<=j and \(\mathrm{Sorted}_W\)(j::v) give i<=j and, by its induction
bound j<=each later index and nat transitivity, i<=each later index.
Also a sorted tail remains sorted. Conversely if w is sorted and
every occurrence in w is at least i, i::w is sorted: empty w uses
the singleton branch; nonempty w uses its head occurrence to supply
the required adjacent bound. These explicit two lemmas prove the
merge invariant by the same nested recursion. Empty cases are given
sorted inputs. If leb i j=one, i<=j; the induction result merges
sorted u with sorted j::v. Every resulting occurrence is in u or
j::v by the membership lemma, and is at least i by the preceding
lower-bound lemmas and transitivity. Prefixing i is therefore sorted.
If the code is O, nat trichotomy and neg(i<=j) give j<i, hence
j<=i. The smaller-second induction result merges i::u with sorted v;
all its occurrences are at least j by the identical argument. Prefix
j for the output. Selector equations in these two code branches are
obtained by congruence and Prop transport, never equality reflection.

Fix the declared semiring and valuation. Nested induction proves the
semantic merge equation even without sortedness. Empty/second cases
use the two multiplicative-unit laws. In the true branch its value is
\(\nu i\times_A(V_\nu u\times_A V_\nu(j::v))\) by the first
induction result; associativity gives the original first-word value
times the second-word value. In the false branch its value is
\(\nu j\times_A(V_\nu(i::u)\times_A V_\nu v)\).
Associativity, commutativity exchanging the first two factors, and
associativity back give
\(V_\nu(i::u)\times_A(\nu j\times_A V_\nu v)\).
That is the required product. These are the only two nonempty branches;
the actual multiplicative commutativity law is essential here and has
not been assigned to an arbitrary uninterpreted carrier.
\end{proof}
\end{annexitem}

\begin{annexitem}{annex:p12.har.poly-insertion}{thm:p12-har-poly-insertion}{The complete insertion induction, key invariant and value equation}
\begin{proof}
Row patterns are the typed single \(\mathrm{Rec}_R\) branch. Its coefficient and
word projections are nat and W. At fixed c:nat,w:W, mk has a
nat code, two P arguments and target P. The inner insertion seed is
P; its branch receives r:R,s:P,z:P and applies only the two typed
word codes, \(I_P\), plus and the polynomial constructor. Its target is
P, so every insertion expression has type P independently of a
Good proof. Sum/Scale/Prod will instantiate these actual values, not
the invariant evidence.

Prove the output-key subset by induction on p, with c,w quantified,
using data cases on c. At c=O the outer eqb code computes to one
and insertion returns p; every output key is an input key. For
c=S k the code is O and the recursive result on a tail is exactly
the same nonzero insertion on that tail. At p=epsilon the sole output
key is w. At head row d v and tail s, equality-code true combines
at key w followed by s; larger-new-key true prepends w to the whole
input; otherwise the output keeps v and the recursive tail. In the
last branch its induction subset is either w or an s key. These are
every possible output key. Membership cases are obtained explicitly
from the head/tail occurrence constructors: at a cons they give
equality with the head row or a tail occurrence; coefficient/word
projection congruence transfers a head equality to its key. No
unimplemented key-set map is assumed. Nat plus(S k)d is S(k+d),
so the combination coefficient is nonzero and its mk test returns
the combined row, even if the raw old coefficient is O.

Now strengthen that same induction to Good preservation for sorted w.
The zero-coefficient branch returns the given Good input. In the
nonzero/empty branch, the output singleton has nonzero c, sorted w,
Good empty tail and no tail key. At a Good head row d v, the old
coefficient is nonzero, its word sorted, the tail Good and every tail
key strictly below v. If weq w v is one, its proved specification
gives w=v. The output combined coefficient is the nonzero S(k+d);
word w is sorted; every tail key below v is below w by Prop equality
transport. This gives all Good-cons premises. If equality is false
and wlt v w is one, prepend the new row. The old head key is below
w by that code; every old tail key is below v and hence below w by
lex transitivity. If both tests are false, equality is excluded and
not(v<lex w); strict trichotomy therefore gives w<lex v. The tail
induction supplies Good insertion into the old tail. Its output-key
subset says each new tail key is w or an old tail key, and both are
strictly below v by this comparison or the old Good premise. Thus
the unchanged old head is Good over this entire new tail. These cases
include every ordering/combination/zero branch. An equation about a
comparison code is only sent through \(I_P\) by congruence; Good is a
Prop predicate, so the corresponding output equations transport this
proof without equality reflection or proof-to-data selection.

For the semantic equation induct on arbitrary raw p with c,w quantified,
again case c. At c=O insertion is p; chi0=0A and zero annihilation
make the proposed new-row value zero, and the additive identity proves
the formula. At nonzero c and empty p, insertion is its singleton;
the value is chi c times Vw plus0A, equal to the right side.
At head row d v let H denote its interpreted value and J the tail
value; let C denote chi c times Vw. If the equality test is true,
w=v by its proved specification. Congruence for V transfers that
word equality. The coefficient homomorphism and distributivity give
\[
 (\chi(c+d)\times_A Vw)+_A J
   =((\chi c+_A\chi d)\times_A Vw)+_A J
   =C+_A(H+_A J).
\]
The combined mk coefficient is nonzero as checked above. In the
prepend branch the value is already C+(H+J), exactly the target.
In the recursive branch the induction hypothesis gives tail value
C+J, so the result H+(C+J) becomes C+(H+J) by addition associativity,
commutativity of H,C and associativity back. This proves the equation
on all raw polynomials without a sortedness premise.

Finally let c be positive and w larger than every input key. Empty
p gives the singleton. With a head key v, the premise gives v<w,
so w differs from v by asymmetry/irreflexivity. The weq code is O
and the wlt(v,w) code one; their constructor-selector equations prove
Ins(c,w,p)=row c w \# p in primitive equality. This is the precise
already-in-position property required by quotation retraction.
\end{proof}
\end{annexitem}

\begin{annexitem}{annex:p12.har.poly-operations}{thm:p12-har-poly-operations}{Every sum, scale/product branch and zero behavior}
\begin{proof}
The body gives three \(\mathrm{Rec}_P\) instances with target P. Each branch
receives its original row and tail plus recursive P result, pattern-
matches the row by its typed \(\mathrm{Rec}_R\) case, and applies Ins, MW or
Sum at their declared data types. This verifies all applications and
abstractions, including both loops in multiplication.

For Good Sum, induct on its first input, with a Good second input
as its seed premise. Empty returns that Good seed. Each nonempty
branch inserts its original row into the already Good recursive result;
its word is sorted by the input Good projection. Insertion accepts
every coefficient, so no hidden positive-coefficient theorem is needed
to form its result. For Scale, the seed is Good epsilon. A source
row's word is sorted; MW merges it with the sorted scaling word and
is sorted by the word theorem. Insert the computed product coefficient
into the Good recursive value. Its possible zero is dealt with by
Ins itself, not ruled out as a missing premise. For Prod, each source
head word is sorted; Scale on the Good second input is Good, and
the recursive tail product is Good. Sum of those gives the stated
Good output. These are all constructor branches.

For value of Sum induct on the first raw input. Empty uses zero plus
the seed value. At row c w \# s the insertion equation and induction
give B(row c w)+(L s+L q); associativity makes (B(row c w)+L s)+L q,
the required value. No data-order premise is used in this equation.
For Scale, fix its c,w and put K=chi c times Vw. At empty input,
both sides are zero by annihilation. At row d v \# s, insertion gives
the new-row value plus recursive value. The coefficient and merge
theorems give its new-row value as
\[
 (\chi c\times_A\chi d)\times_A(Vw\times_A Vv)
       = (\chi c\times_A Vw)\times_A(\chi d\times_A Vv).
\]
Explicitly reassociate the four factors, commute the middle chi d
and Vw, then reassociate; these are the declared multiplicative laws.
The induction value is K times L s, so distributivity combines the
sum into K times(B(row d v)+L s), the asserted scale equation.
For Prod, at empty input zero annihilates L q. At row r \# s, the
Sum equation, Scale equation and induction give
(B r times L q)+(L s times L q). Right distributivity yields
(B r+L s)times L q, exactly the product interpretation. This covers
every multiplication/combination case and explicitly uses commutativity
where factor permutations occur.

The zero behaviors are actual structural recursions too. Sum(empty,q)
is its seed q. Prod(empty,q) is its empty seed. Scale(c,w,empty) is
empty. At c=O, each scale coefficient O times d contracts to O
and Ins(0,word,z) returns z, so P induction gives Scale(0,w,p)=empty.
With second input empty, every Prod head scale is empty; Sum(empty,z)
returns z, so its P induction gives Prod(p,empty)=empty. For literal
constructor data these equalities give finite constructor reductions;
for arbitrary open data their displayed structural proofs are book
equalities, not an assertion that an opaque input has already reduced.
\end{proof}
\end{annexitem}

\begin{annexitem}{annex:p12.har.poly-normalization}{thm:p12-har-poly-normalization}{Every original grammar constructor and the universal data invariant}
\begin{proof}
The source T fold here has target P. Its zero seed is P epsilon;
its unit seed is a P singleton; its variable branch has type
nat-to-P. Each binary branch has original s,t:T and recursive values
p,q:P, then returns the typed Sum or Prod. Thus N:T-to-P is the
complete actual five-branch data program, with no normal-form proof
as an input.

Induct on T into the conjunction of Good(Nt) and the interpretation
equation, with the fixed semiring and arbitrary typed valuation as
parameters. At Z, N returns epsilon: Good's empty constructor applies
and L epsilon=0A=E Z. At U, it returns row one [] \# epsilon;
one is separated from O, the empty word sorted and its tail empty.
Its value is chi(one)times1A+0A=1A by the proved coefficient and
semiring unit/zero laws, which is E U. At Var i it returns row one
(i::[]) \# epsilon: singleton sortedness supplies its word premise,
and its value is 1A times(nu i times1A)+0A=nu i, the actual variable
interpretation. These prove all three initial branches.
At Add s t, both induction values are Good and equal to the two
source interpretations. Good Sum proves the output invariant; its
value equation and two congruences/transitivity give
L(Sum(Ns,Nt))=E s+E t=E(Add s t). At Mul s t, use Good Prod,
its value equation and the two induction congruences in the identical
way with multiplication. This yields L(Prod(Ns,Nt))=E(Mul s t).
Those are the two remaining constructors. All valuations were arbitrary;
none were restricted to closed numerical inputs or assumed extensionally
equal. The only algebraic premises are the explicitly declared semiring
equations used in the helper proofs.
\end{proof}
\end{annexitem}

\begin{annexitem}{annex:p12.har.poly-quotation}{thm:p12-har-poly-quotation}{Every literal/word/row/polynomial quotation case and its value}
\begin{proof}
All four quotation targets are T. Lit's nat step receives k:nat and
t:T then applies the typed \(I_T\) to U or Add t U. The word step
receives i:nat,w:W,t:T and selects Var i or Mul(Var i)t. The
single row step receives c:nat,w:W and selects one of three T
expressions. The polynomial step receives r:R,p:P,t:T and selects
\(Q_R\) r or Add(\(Q_R\) r)t. Thus every branch has the precise generated
recursor type and each displayed quotation is actual data, not the
existence of an unspecified expression with the desired value.

For lit use nat induction. At O the quotation is Z with value 0A,
which is chi0. At S k, case k. For k=O its eqb code is one and
the quotation contracts to U; chi(one)=1A proves the goal. For
k=S j the eqb code is O and the quotation is Add(lit k)U. Its
value is chi k+1A by the induction hypothesis, precisely chi(S k)
by its defining computation. These cases include every coefficient.

For \(Q_W\) use word induction and case its tail. At [] the quotation
U has value1A. At i::[] its quotation Var i has value nu i, whereas
the word interpretation is nu i times1A; the unit law proves equality.
At i::(j::v) the quotation is Mul(Var i)(\(Q_W\)(j::v)). The recursive
equality gives nu i times V(j::v), exactly the cons word value.
The empty test is a data constructor case, with no omitted neutral
or proof-selected input.

For \(Q_R\) pattern its sole constructor c,w and case w. At empty w,
the quotation lit c has value chi c, and its row value is chi c
times1A; use the unit law. At nonempty w split the actual eqb(c,one)
range. In its one branch c=one by the equality-code theorem. The
selector-output equation is obtained by congruence through \(I_T\) and
its one computation. The \(Q_W\) value and chi(one)=1A then equal the
row value using the multiplicative unit law. In its O branch the
selector gives Mul(lit c)(\(Q_W\) w), whose two already proved values
give chi c times Vw by constructor computation and two congruences.
This covers c=O too: the quoted product has a zero literal and its
value is zero by annihilation, as is the row interpretation.

For \(Q_P\) use P induction, then case the tail. At epsilon, Z has value
0A. At a singleton r\#epsilon, the quotation is \(Q_R\) r, whose proved
value B r equals B r+0A=L(r\#epsilon) by additive identity. At r\#(s\#p)
the quotation is Add(\(Q_R\) r)(\(Q_P\)(s\#p)); the row equation and recursive
tail equation give B r+L(s\#p), exactly its list interpretation.
These are all polynomial shapes. Code equations about an open c were
used only via congruence followed by Prop equality transport, never
as definitional conversion. No semiring-model completeness assertion
or equality of interpretation functions is needed.
\end{proof}
\end{annexitem}

\begin{annexitem}{annex:p12.har.poly-preservation}{thm:p12-har-poly-preservation}{The full universal open-valuation chain and its Nat instance}
\begin{proof}
Fix an arbitrary valid context with the declared A:Set, values,
operations, their commutative-semiring proof families and nu:nat-to-A.
All operations G,N,\(Q_P\),\(E_\nu\) and \(L_\nu\) have already been explicitly
typed and their helper proofs quantify over these same parameters.
For every t:T, the actual definition gives Gt=\(Q_P\)(Nt). The quotation
theorem instantiated at this data Nt supplies
\(E_\nu^A(Q_P(Nt))=L_\nu(Nt)\).
The complete five-constructor normalization induction supplies
\(L_\nu(Nt)=E_\nu^A t\). Compose these two primitive equalities
and expand the finite G abbreviation. This proves exactly the universal
equation, with no semantic-preservation premise assumed in a helper.

The source Nat instance takes A=nat, O,one, plus,times. Additive
associativity/commutativity/zero and multiplicative associativity/
commutativity/unit/zero/distributivity have the supplied historical
arithmetic proofs in thm:p12-har-prime-semiring and its annex, including
right distribution derived from left distribution and commutativity.
Thus the actual model premises are discharged there; no DN is used.
Since nu was arbitrary, this includes every well-typed open natural
valuation and the explicitly constructed finite-context valuations,
not just closed inputs. No function-extensionality theorem converts
pointwise assignments into equality of functions. Conversely nothing
in the proof equips an arbitrary operation-only A with commutativity;
those laws are genuinely required by word permutation and product
collection and remain explicit for every other interpretation.
\end{proof}
\end{annexitem}

\begin{annexitem}{annex:p12.har.poly-retraction}{thm:p12-har-poly-retraction}{Constructed atomic normalizations, normal-form retraction and idempotence}
\begin{proof}
First supply the data identities needed by retraction, instead of
assuming semantic preservation implies equality of sparse records.
For all c,w,p, Sum(single(c,w),p)=Ins(c,w,p): at c=O both single
and insertion are empty/unchanged by their zero computation; at c=S k
the Sum singleton fold has recursive seed p and inserts its one head.
Also
\[
 \mathsf{Prod}(\mathsf{single}(c,u),\mathsf{single}(d,v))
       =\mathsf{single}(cd,\mathsf{MW}(u,v)).
\]
If c=O both sides are empty by source times zero and Prod's zero
case. If d=O, right Prod zero and timesaltzero give the same conclusion.
For successor c,d, expand the one-row product: its Scale inserts cd
and MW u v into empty, producing the right-hand single; the outer
Sum with empty leaves this singleton (including any zero through the
single/Ins zero identity). These expansions prove every coefficient
case, not a formal-semiring completeness axiom on P.

For constants,
\(\mathsf{Sum}(\mathsf{single}(c,[]),\mathsf{single}(d,[]))
=\mathsf{single}(c+d,[])\).
If c=O the Sum seed and plus zero give it. If c is a successor and
d=O, empty insertion and \(\mathsf{plus\_altzero}\) give it. If both are successors,
weq([],[]) computes one, so insertion combines them at the empty
word; their sum is a successor, so mk keeps exactly that single.
Nat induction now proves N(lit c)=single(c,[]): zero is N Z;
the first successor quotes U; later successors quote Add(lit k)U,
whose N is this proved constant sum with coefficient one, namely k+one.
The source identity k+one=S k gives the required data equality.

Next, if \(\mathrm{Sorted}_W\)(i::u), then MW(i::[],u)=i::u. At empty u it
is the literal empty-second merge case. At j::v its sortedness gives
i<=j, hence the leb code one by the actual comparison theorem.
Congruence through the merge selector, then MW([],j::v)=j::v,
gives exactly i::j::v. This is a book equality for opaque indices,
not an assertion that a comparison on those indices already contracted.
Word induction using this identity proves
N(\(Q_W\) w)=single(one,w) for sorted w. Empty quotes U. A singleton
quotes Var i. For a longer i::u, the quote is Mul(Var i)(\(Q_W\) u);
the tail induction and the proved singleton product give coefficient
one-times-one=one and word MW(i::[],u)=i::u. Both data fields
are restored by constructor congruence.

For a positive coefficient c and sorted w, prove
N(\(Q_R\)(row c w))=single(c,w). Empty w is the literal result already
proved. For nonempty w split eqb(c,one). Its one branch gives c=one
and quotes \(Q_W\) w, so the word normalization just proved and coefficient
transport give the statement. Its zero branch quotes Mul(lit c)(\(Q_W\) w);
the two atomic normalizations and singleton product give coefficient
c times one=c and word MW([],w)=w. These are every row quotation
branch. The input positivity assures single(c,w) is its actual row,
not an unretained zero case.

Now use P induction with Good p as premise. At empty p, Q_Pp=Z
and N returns empty. At row c w \# s, Good supplies positivity,
word sortedness, Good s and strict domination of every s key.
If s is empty, \(Q_P\) returns the row quotation and its atomic result
is exactly p. If s is nonempty, \(Q_P\) returns Add(\(Q_R\) row)(\(Q_P\) s).
The row result and induction result give
Sum(single(c,w),s)=Ins(c,w,s). The already-in-position insertion
lemma applies to the supplied strict-domination premise and returns
row c w \# s, precisely p. These two tail cases close the retraction
without a claim that equal evaluations in one model imply equal data.

For arbitrary t, Good(Nt) is already proved. Instantiate retraction
there: N(\(Q_P\)(Nt))=Nt. Congruence for \(Q_P\) and expansion of the actual
G definition give G(Gt)=Gt. The representation witness for its normal
form is the already computed data Nt, with its actual Good proof;
no data is chosen from a Prop existential. If desired, equality of
two quotations of Good records implies equality of their records
by congruence for N and these retraction equations. None of this
asserts model-semantic completeness in a semiring where distinct
natural coefficients or variables may have equal values.
\end{proof}
\end{annexitem}

\begin{annexitem}{annex:p12.har.poly-example}{thm:p12-har-poly-example}{The actual source-pattern computation, zero cases and association qualification}
\begin{proof}
Let i be the literal O name index, X=Var O, w=[O], and ww=[O,O].
Their word fields are constructor data, so every required comparison
really contracts: eqb O O is one, leb O O is one, weq is its
cons/empty recursion, and []<lex w<lex ww has the corresponding
literal wlt one calculations. The numeral coefficient one-times-one
contracts to one and one+one to source two.

N X is single(one,w), and N U is single(one,[]). The Sum fold
inserting w into the unit constant compares []<lex w, so both N
copies of Add X U are the actual descending polynomial
\([\mathsf{row}\ \mathsf{one}\ w,
\mathsf{row}\ \mathsf{one}\ []]\).
Here brackets abbreviate the explicitly declared \#/epsilon constructors.
Scale of its variable row into that same polynomial first has the
tail product word MW(w,[])=w and then the head product MW(w,w)=ww.
Insertion places the ww row before w, yielding [row one ww,row one w].
Scale of its constant row produces [row one w,row one []], since
MW([],v)=v. The Prod Sum fold first inserts the former w row into
that latter polynomial; weq(w,w) is one, so coefficients combine
to two at w. It then inserts ww in front, using wlt(w,ww)=one.
The final computed sparse data are exactly
\([\mathsf{row}\ \mathsf{one}\ ww,
\mathsf{row}\ \mathsf{two}\ w,
\mathsf{row}\ \mathsf{one}\ []]\).
This explicitly covers the monomial ordering, repeated-variable merge
and like-term combination; it is not inferred from a polynomial identity.

Quotation of ww with coefficient one is Mul X X: the terminal
word index quotes Var O rather than Var O times U. At the coefficient
two row, eqb two one is O, lit two computes Add U U, and the
nonempty word w quotes X; its row is Mul(Add U U)X. The last empty
word row quotes lit one=U. \(Q_P\) therefore returns the exact right-
associated AST in the theorem, with no trailing zero. The constructor
steps just listed give a finite beta/iota/delta reduction path for G
on this literal quotation; no external evaluator was executed.

Under any declared model valuation the actual output reads
x-times-x+((1A+1A)-times-x+1A). Addition associativity is the explicit
equation identifying this interpretation with the left-associated
(x-times-x+(1A+1A)-times-x)+1A reading. The source's exponent/two
notation is precisely these arithmetic printing abbreviations. Thus
the source example's polynomial is obtained with the declared
parenthesization convention, without falsely claiming literal equality
to an alternate AST parsing. In the Nat instance this is \(x^2\)+(2x+1)
and its source association bridge.

For clarity the zero cases are also genuine computations: N Z is
empty and \(Q_P\) empty is Z; Prod(empty,q) discards its right argument;
right-empty product and zero scaling structurally return empty by
the proved zero clauses; Ins(0,w,p) returns p before inspecting p.
Constants use the empty word, repeated powers use retained repeated
indices, and the coefficient-one test omits a spurious unit factor.
No source grammar constructor or ordering/zero case was delegated to
an unimplemented normalization or library routine.
\end{proof}
\end{annexitem}
